\chapter{Continual learning and population adaptation}
\label{evod50:chapter}
Train a two-coordinate learner toward $(1,0)$, then toward $(0,1)$.
Should it finish at the new target, the midpoint, or somewhere else? Each
answer optimizes a different objective. A memory store that preserves old
examples, a regularizer that resists movement and a population search that
chooses among learners introduce different costs and causal explanations.
We derive those differences before using the word adaptation as evidence.

A quadratic learner and an exploit/mutate example form the numerical spine. We add
an independent closed-form oracle, explicit retention conventions,
replay gradients, feasibility constraints and counted search opportunities.
All reported values are synthetic arithmetic. No DOGMA or Hermon training
improvement follows from them.

\section{Five adaptation layers and their owners}
\index{parameter learning}\index{structural adaptation}\index{inheritance}
\index{continual learning!adaptation layers}
Context updates change current inputs. Durable memory updates change later
retrieval opportunities. Parameter or adapter updates change the function.
Structural changes alter its graph, routing or memory topology. Population
updates change which candidates exist and are selected. A saved conversation
changes neither weights nor architecture merely by being stored.

Let $h_t$ be inference state, $M_v$ a memory version, $\theta_k$ weights,
$G_s$ a structure and $\mathcal P_g$ a population. The indices name update
events, not a universal ordering of time scales. A deployment should identify
which versions a request sees. Changing $\theta$ while retaining an incompatible
native carry or K/V state can silently change the request's meaning.

Chapter 13's search inherited masks with fresh weights; its alternative
inheritance map was separately defined. The same discipline applies here.
Copying weights alone differs from copying weights, optimizer moments,
normalization statistics, data cursor and RNG. A structural mutation needs
a semantic coordinate map, not merely tensors of matching length. The local
reference is not connected to Chapter 43's governance registry; it neither
self-modifies a runtime nor promotes an artifact.

\section{A transparent acquisition--retention conflict}
\index{acquisition}\index{catastrophic forgetting}
Let $\theta\in\mathbb R^2$ and task target $u\in\mathbb R^2$ be dimensionless.
Define a task loss and a bounded teaching score
\[
L_u(\theta)=\tfrac12\|\theta-u\|_2^2,
\qquad Q_u(\theta)=\max\{0,1-\|\theta-u\|_2^2/4\}.
\]
$Q$ is not classification accuracy; its scale is chosen to make retention
visible. The reference function is still named \texttt{accuracy\_matrix}; it
returns this score, not accuracy. Clipping can hide
differences between very poor candidates, so an actual study should retain
the underlying loss as well.

Gradient descent with rate $\eta$ gives
$\theta_{k+1}=(1-\eta)\theta_k+\eta u$.
Subtracting $u$ and iterating yields
\[
\theta_n-u=(1-\eta)^n(\theta_0-u).
\]
For $0<\eta<2$ the error contracts. Starting at zero, four steps at $\eta=1/2$
toward $u_1=(1,0)$ give $\theta_A=(15/16,0)$. Four more toward $u_2=(0,1)$ give
\[
\theta_B=(15/256,15/16).
\]
The new task is almost fitted; the first coordinate has mostly been forgotten.
This is the objective's behavior, not a software bug.

\Plate{01-quadratic}{Parameter space contains two task optima. Four half-rate
updates approach task one, then task two. An anchor changes the second objective
and converges to a compromise. The path is a deterministic numerical fixture,
not a trajectory measured in a trained neural network.}{evod50:quadratic}

\section{Derive an anchored objective and its stability bound}
\index{elastic weight consolidation}\index{learning rate!curvature bound}\index{closed-form oracle}
EWC uses a quadratic penalty weighted by parameter importance estimates
\cite{evod50:ewc}. In the usual approximation, a diagonal Fisher-information
estimate represents local importance around a previous optimum. The original
paper discusses approximations and uncertainty limits; the method does not
guarantee retention on arbitrary new tasks. Our teaching fixture sets
$F_i=1$ explicitly. It does not estimate a Fisher matrix and is therefore
an EWC-like anchored quadratic, not a reproduction of the complete method.

With anchor $a$ and nonnegative weights $F_i$, let
\[
L(\theta)=\tfrac12\|\theta-u\|^2+
\tfrac\lambda2\sum_i F_i(\theta_i-a_i)^2,\qquad\lambda\ge0.
\]
Coordinate differentiation gives
\[
g_i=(\theta_i-u_i)+\lambda F_i(\theta_i-a_i),\quad
\theta_i^*=\frac{u_i+\lambda F_i a_i}{1+\lambda F_i}.
\]
Set $c_i=1+\lambda F_i$. Subtracting the optimum from one gradient step gives
\begin{equation}\label{evod50:contraction}
\theta_{k+1,i}-\theta_i^*=(1-\eta c_i)(\theta_{k,i}-\theta_i^*),
\quad
\theta_{n,i}=\theta_i^*+(1-\eta c_i)^n(\theta_{0,i}-\theta_i^*).
\end{equation}
Thus all coordinates contract when
$0<\eta<2/\max_i c_i$. Stronger anchoring increases curvature and can require
a smaller learning rate. An implementation can be arithmetically correct while
its optimization diverges. The negative control takes $\eta=3$ on the
unanchored scalar loss and observes increasing error.

\Listing{step}
\Listing{closed_form}
With $a=\theta_A$, $F=(1,1)$, $\lambda=1$ and $\eta=1/2$, the contraction
factor is zero. One second-task step reaches $(15/32,1/2)$ exactly, and the
remaining steps leave it there. Old-task retention improves relative to the
naive path, while new-task acquisition worsens. This tradeoff is the central
result of the fixture; reporting only reduced forgetting would hide its cost.

\begin{center}\input{\Generated/adaptation.tex}\end{center}
The generated values use four updates per experience and the stated synthetic
score. The tests compare loops with the affine formula, exact rational first-task
values and a finite-difference gradient. These checks have different failure
modes: a wrong update count, derivative or normalization can be detected without
reusing the same iterative implementation as its own oracle.

\section{Read the full experience matrix}
\index{forgetting!peak convention}\index{backward transfer}\index{forward transfer}
For experiences $E_1,\ldots,E_T$, let $A_{t,j}$ be the evaluation score on task
$j$ after training through task $t$. Row $t$ uses one fixed artifact; column
$j$ uses one fixed evaluation distribution. The diagonal measures immediate
acquisition, and entries below it measure later retention.

\Plate{02-matrix}{A matrix entry binds one model time and one evaluation task.
Diagonal cells mark acquisition; lower cells are retention; upper cells concern
pre-acquisition behavior and are not training data. Future evaluations must stay
hidden from the learner and selection process.}{evod50:matrix}

Average final score is $T^{-1}\sum_jA_{T,j}$. We define nonnegative peak
forgetting for each task as
\[
F_j=\max_{j\le t\le T}A_{t,j}-A_{T,j}.
\]
Including the final row makes this quantity nonnegative. Taking the maximum over
all earlier rows instead would include pre-acquisition values; the post-acquisition
convention used here avoids that.
Last-task forgetting is zero by definition, not evidence of long-term retention.
\Listing{forgetting}
Signed backward transfer is
\[
\mathrm{BWT}=\frac1{T-1}\sum_{j<T}(A_{T,j}-A_{j,j}),\qquad T>1.
\]
It can be positive even when peak forgetting is zero. For $A_{1,1}=.6$ and
$A_{2,1}=.8$, backward transfer is $.2$, while peak forgetting is zero.
Forward transfer needs a target-only or pretraining baseline before acquisition;
an off-diagonal number alone does not supply that counterfactual.

The two-task oracle allows the analyst to calculate both task scores at every
row, since both quadratics are known. The learner's update uses only the current
target and its permitted anchor. In a real online protocol, future task labels
cannot guide updates or selection. Declare whether task IDs and boundaries are
given, how often evaluations occur, and which evaluations may be inspected.
Avalanche is prior work on a shared continual-learning experimental framework
\cite{evod50:avalanche}; using a common framework does not itself make two
data-access protocols equivalent.

\section{Replay and modularity change different terms}
\index{replay!mixture objective}\index{adapters}
If a replay batch mixes old and new losses with old-data fraction $\rho$,
\[
L_{\mathrm{mix}}=(1-\rho)L_{u_2}+\rho L_{u_1},\qquad
\nabla L_{\mathrm{mix}}=\theta-[(1-\rho)u_2+\rho u_1].
\]
The fixture therefore uses the mixture target directly. At $\rho=1/2$,
starting from zero with half-rate learning produces $(.25,.25)$ in one step.
Replay with a finite biased buffer does not equal the exact old-data objective.
Report buffer bytes, selection policy, old/new exposures and privacy restrictions.
\Listing{replay_step}
Adapters can isolate new parameters but still share routing, normalization or
base representations. A router regression can prevent access to an intact
old adapter. A no-interference claim therefore needs end-to-end task tests,
not merely a check that old weight bytes were unchanged. A joint/offline
reference has access to all allowed data and a different training opportunity;
call it an upper reference only under the actual task/optimization assumptions,
not a guaranteed mathematical upper bound on every algorithm.

\section{Population search counts all its opportunities}
\index{population search}\index{budget parity}
Population-based training interleaves model training with exploitation and
hyperparameter exploration \cite{evod50:pbt}. Our preserved fixture is much
narrower: evaluate each two-coordinate candidate on each target, retain the
best mean quality, then create two deterministic coordinate mutations. It has
no asynchronous workers or learned hyperparameter schedule. Ties retain the
first candidate in declared input order, so completion speed cannot decide them.

For $N$ candidates and $J$ tasks the selection uses $NJ$ candidate-task
evaluations. A budget below $NJ$ rejects before selection. This count excludes
candidate training and runtime overhead, which actual studies must also charge.
Offspring are proposals and require their own future evaluation. Their parent
score cannot certify the changed coordinates.

\Plate{03-population}{Three candidates receive two task evaluations each,
consuming six opportunities. The selected parent supplies inherited coordinates
to two new proposals. Dashed descent indicates inheritance; it does not carry
the parent's evaluation evidence to the mutations.}{evod50:population}

Selection only by current mean can erase complementary skills. Diversity can
be measured in behavior, architecture or failures; parameter distance need not
imply behavioral diversity. A useful evaluation freezes population size,
mutation rate, inheritance map, training budget and selection data. Compare
equal total work as well as per-candidate work, and reserve a final evaluation
that search never queries. Repeated search on a public test is development.

\section{Feasibility, Pareto comparisons and poisoning}
\index{Pareto dominance}\index{poisoning}\index{rollback}
A scalar fitness $q-\alpha c-\beta r-\gamma f$ combines unlike units only
after explicit scaling and policy weights. A risk constraint is often better
expressed as a hard admission condition. Our reference first removes candidates
above a stated risk ceiling, then computes the Pareto set for quality
(maximize), cost and forgetting (minimize). Candidate A dominates B if it is
no worse on every axis and strictly better on at least one. Incomparable
candidates remain; the frontier is not an automatic deployment choice.
\Listing{pareto}
The synthetic candidates \texttt{a} and \texttt{b} survive, a dominated
candidate is removed, and a high-quality unsafe candidate fails feasibility.
Risk here is a supplied number, not a measured probability or security
certification. Tests enforce the mathematical selection rule only.

For a real update, preserve parent lineage, data admission decisions, optimizer,
seeds, artifact identity, acquisition/retention matrix, calibration, budgets and
failure cases. Quarantine untrusted experience before training: poisoning can
persist in memory and weights. An anomaly filter has false negatives; no one
filter guarantees clean training. Ablate replay, importance penalty, mutation
and selection separately, including matched random-change controls.

Promotion remains an external authority decision. A canary needs bounded scope,
duration, observable failure thresholds and a rollback artifact. Rollback restores
future behavior under an old version; it cannot erase messages already emitted
or outside actions already performed. Correct data deletion may require removing
derived memory and retraining artifacts, not merely changing a version pointer.

\section{Selective glossary}
\begin{description}
\item[Acquisition.] Performance on a task after its learning experience under
the declared update budget.
\item[Peak forgetting.] Best post-acquisition score minus final score, including
the final row in the maximum under this chapter's convention.
\item[Backward transfer.] Signed final-minus-acquisition score on earlier tasks;
positive values indicate improvement rather than loss.
\item[Replay mixture.] A training objective combining permitted old and new
data, with memory and exposure costs explicitly counted.
\item[Pareto dominance.] No worse on every chosen objective and strictly better
on at least one, after hard feasibility constraints are applied.
\end{description}

\section{Exercises and worked solutions}
\begin{enumerate}
\item Saving a dialogue changes which layer? \textbf{Answer:} durable evidence
or context, unless a separate parameter update actually runs.
\item Derive unanchored $\theta_n$. \textbf{Answer:} subtract $u$ from the
update and iterate the factor $(1-\eta)$.
\item Four half-rate steps from zero toward $(1,0)$? \textbf{Answer:}
$(15/16,0)$.
\item Four further steps toward $(0,1)$? \textbf{Answer:} $(15/256,15/16)$.
\item Is $Q$ classification accuracy? \textbf{Answer:} no; it is an explicitly
clipped quadratic score and can hide differences below its floor.
\item Derive anchored optimum. \textbf{Answer:} set the coordinate gradient
to zero, yielding $(u_i+\lambda F_i a_i)/(1+\lambda F_i)$.
\item Why may stronger anchoring destabilize an old rate? \textbf{Answer:}
curvature becomes $1+\lambda F_i$, tightening the contraction bound.
\item Does $F_i=1$ reproduce Fisher estimation? \textbf{Answer:} no; it is
a supplied identity importance matrix for a quadratic fixture.
\item What is the one-step optimum for the anchored second task? \textbf{Answer:}
$(15/32,1/2)$ under the stated half-rate and unit penalty.
\item Why report both task scores? \textbf{Answer:} retention can improve by
refusing new acquisition; reduced forgetting alone hides this cost.
\item Last-task peak forgetting? \textbf{Answer:} zero by definition, since
only its acquisition/final row has been observed.
\item Old score improves $.6\to.8$. \textbf{Answer:} BWT $.2$, nonnegative
peak-forgetting convention zero.
\item Does an upper matrix cell prove forward transfer? \textbf{Answer:}
no; compare the permitted pre-acquisition behavior with a target-only baseline.
\item Derive the replay target. \textbf{Answer:} combine gradients linearly,
giving $(1-\rho)u_2+\rho u_1$.
\item Three candidates, two tasks, budget five? \textbf{Answer:} reject;
six candidate-task evaluations are needed before any new offspring evaluations.
\item Does a parent's score validate its mutation? \textbf{Answer:} no;
changed coordinates define a new candidate.
\item Why can two Pareto survivors remain? \textbf{Answer:} one can have
higher quality and higher cost, so neither dominates the other.
\item Design an adaptation ablation. \textbf{Rubric:} frozen/sequential/replay/
anchor controls, declared task access, equal total data and search budget,
seed distributions, full retention matrix, calibration and untouched final tests.
\item Design a poisoning test. \textbf{Rubric:} bounded injected experience,
source provenance, quarantine controls, targeted and clean-task evaluations,
persistence/deletion checks and no claim of security certification.
\item Can rollback undo outside side effects? \textbf{Answer:} no; restore
artifact behavior separately from reconciling external state.
\end{enumerate}

\section{Bridge}
Adaptation mechanisms can pass every numerical contract and still fail to
acquire useful skills. Chapter 51 turns that gap into operational capability
hypotheses with explicit evidence and falsifiers.
